\section{Joining transfers through planetary encounters}
\label{sec:encounters}
\subsection{An unpowered point event}
In the zero-duration patched-conic model, let the planet state at
encounter epoch $\tau$ be $(R_p,V_p)$ with acceleration $A_p$.
Use common heliocentric units, incoming excess velocity
$w_-=v^--V_p$, speed $s=|w_-|>0$, and rotation $Q(\alpha)$.
The event map is
\begin{equation}
 r^-=r^+=R_p,\quad
 v^+=V_p+Q(\alpha)w_-,\quad b^+=b^-.
 \label{eq:pointflyby}
\end{equation}
The position is a mathematical encounter node, not a physical collision.
With normalized planetary parameter $\bar\mu_p$ and minimum
planetocentric periapsis $q_{\min}$,
\begin{equation}
 |\alpha|\le D(s)=2\arcsin(1/e),\qquad
 e=1+\frac{q_{\min}s^2}{\bar\mu_p}.
 \label{eq:turn}
\end{equation}
This preserves excess speed and consumes no propellant.
The heliocentric energy change is
$\Delta E_\odot=V_p\cdot(Qw_--w_-)$, and
\begin{equation}
 |\Delta v|\le\frac{2s}{1+q_{\min}s^2/\bar\mu_p}
 \le\sqrt{\bar\mu_p/q_{\min}}.
\end{equation}
Zero excess speed is not a hyperbolic flyby. Zero turn is an
infinite-periapsis limit of this event set; a finite-SOI model treats
the geometry explicitly.

A circular-rendezvous atlas fixes the ship velocity to the planet's,
so it has $w_-=0$ at that endpoint and cannot supply general flyby legs.
General endpoint actions are needed. Once they exist, one encounter
can be parameterized by $(w_{-,x},w_{-,y},\alpha,\tau)$.
For the sum $F$ of incoming and outgoing actions, let superscripts
$-,+$ denote their adjacent endpoint covectors and put
$w_+=Qw_-$. Equation~\eqref{eq:endpoint} gives
\begin{equation}
 \begin{split}
 F_{w_-}&=p_v^- -Q^Tp_v^+,\\
 F_\alpha&=-p_v^+\cdot\rot w_+,\\
 F_\tau&=H^+-H^-+(p_r^- -p_r^+)\cdot V_p\\
 &\hspace{18mm}+(p_v^- -p_v^+)\cdot A_p.
 \end{split}
 \label{eq:pointgrad}
\end{equation}
Here the planet-relative $w_-$ is held fixed when varying epoch;
both endpoint velocities acquire $A_p\,d\tau$.
Specifically, the incoming endpoint differential is
$p_r^-\cdot V_p\,d\tau+p_v^-\cdot A_p\,d\tau-H^-d\tau$;
the outgoing initial differential has the opposite covector signs and
$+H^+d\tau$. Their sum is the stated $F_\tau$.

Let $\eta_\pm\ge0$ multiply $\pm\alpha-D(w_-)\le0$ and
$d_D=\nabla D=-4(q_{\min}/\bar\mu_p)w_-/
(e\sqrt{e^2-1})$. Interior-event stationarity requires
\begin{equation}
 \begin{gathered}
 p_v^-=Q^Tp_v^++(\eta_++\eta_-)d_D,\\
 p_v^+\cdot\rot w_+=\eta_+-\eta_-,\qquad F_\tau=0,
 \end{gathered}
\end{equation}
together with complementary slackness. An interior turn need not be
maximum turn. At a prescribed periapsis and turn side $\sigma=\pm1$,
the velocity Jacobian is
\begin{equation}
 M=Q+(\rot w_+)(\sigma d_D)^T,\qquad p_v^-=M^Tp_v^+.
\end{equation}
The extra term differentiates the speed-dependent turn; $\det M=1$.
Resource costates match if the same resource passes unchanged.
The position guard supplies an independent position-covector multiplier.
The matched Hamiltonian combination is
$H-p_r\cdot V_p-p_v\cdot A_p$, rather than $H$ alone.

\subsection{A finite powered encounter}
For an encounter with duration, use a planet-centered nonrotating patch,
length scale $L$ and local units
$T_p=\sqrt{L^3/\mu_p}$, $V_p^*=\sqrt{\mu_p/L}$,
$J_p=(V_p^*)^2/T_p$.
The controlled equations again have form \eqref{eq:kepler}, with
$x$ the relative position and
\begin{equation}
 \varepsilon\le|x|\le\sigma_S,\quad
 \varepsilon=r_{p,\min}/L,\quad \sigma_S=R_S/L.
\end{equation}
Choose inward entry and the first outward exit. With $L=R_S$,
$\sigma_S=1$. Distinct loop branches require explicit labels.
The engine selector and radial-floor conditions are those of
Section~\ref{sec:constraints} in the new units.

At an endpoint with positive planetary energy
$E=|v|^2/2-1/|x|$, put $s=\sqrt{2E}$, $j=x\times v$ and
\begin{equation}
 e_v=(|v|^2-1/|x|)x-(x\cdot v)v,\quad e=|e_v|,\quad n=e_v/e.
\end{equation}
The hypothetical ballistic tails have excess-velocity labels
\begin{equation}
 w_{\infty,-}=\frac{s}{e}(n+js\rot n),\qquad
 w_{\infty,+}=\frac{s}{e}(-n+js\rot n).
\end{equation}
Use the incoming expression at entry and outgoing expression at exit.
Thrust can change their magnitudes. The unpowered equal-speed and
maximum-turn constraints must not be imposed on this controlled phase.
An $E\le0$ endpoint has no hyperbolic excess label, although finite-SOI
crossing can still occur.

The exact work and angular-momentum identities are
\begin{equation}
 E'=v\cdot u,\qquad j'=x\times u.
 \label{eq:work}
\end{equation}
For positive-energy labels,
$(s_+^2-s_-^2)/2=\int v\cdot u\,dt$.
High periapsis speed increases instantaneous work per thrust increment
when thrust is suitably aligned. The mission optimum must also satisfy
exit geometry, timing and resource conditions. It does not follow by
maximizing this work alone.
An impulse is not a constant-power burn: compressing a nonzero
$\Delta v$ into duration $\delta$ costs at least $|\Delta v|^2/\delta$.
Likewise infinity-to-infinity hyperbolic flight has infinite duration;
the optimized mission segment has finite SOI endpoints.

On a connected unconstrained ideal phase,
$p_v''=Dg(x)^Tp_v$. A positive-length coast forces $p_v=p_v'=0$
there, hence everywhere by uniqueness. A nontrivial powered ideal
extremal therefore has no positive-length coast interval.
The finite-band selector admits coasts without $p_v=0$ and has no such
restriction.

\subsection{The moving SOI}
Return to one common physical or heliocentric unit system.
For ship coordinates $(X,V)$ put $x=X-R_p(t)$ and
$v_{\rm rel}=V-V_p(t)$. The consistent global patch fields are
\begin{equation}
 G_\odot=-\mu_\odot X/|X|^3,\qquad
 G_P=A_p-\mu_p x/|x|^3.
 \label{eq:fields}
\end{equation}
Subtracting the origin acceleration gives the desired pure planetary
relative dynamics. Solar tides inside the patch are omitted.

The SOI guard is $c_S=(|x|^2-R_S^2)/2=0$. At a free transverse
crossing with continuous state and mass, endpoint variation gives
\begin{equation}
 \begin{gathered}
 p_X^+=p_X^-+\zeta x,\quad
 p_V^+=p_V^-,\quad p_b^+=p_b^-,\\
 H^+-H^-=\zeta x\cdot V_p.
 \end{gathered}
 \label{eq:soijump}
\end{equation}
These are hybrid junction conditions \cite{PakniyatCaines2021}.
The equality multiplier $\zeta$ has no sign restriction; it is different
from a radial exclusion multiplier. With unchanged engine and a unique
regular selector,
\begin{equation}
 \zeta=-\frac{p_V\cdot(G^+-G^-)}{x\cdot v_{\rm rel}}.
 \label{eq:zeta}
\end{equation}
To derive this, substitute $\Delta p_X=\zeta x$ into the Hamiltonian
jump: $\Delta H=\zeta x\cdot V+p_V\cdot\Delta G$.
Equating it to $\zeta x\cdot V_p$ and dividing by the nonzero transverse
speed $x\cdot(V-V_p)$ yields \eqref{eq:zeta}.
Grazing invalidates this chart. Coincident constraint events require
their additional multipliers.

A useful optimization chart is $(\theta,v_r,v_t,\tau)$:
$X=R_p+R_Se_r(\theta)$,
$V=V_p+v_re_r+v_te_\theta$.
With $d_r=p_X^- -p_X^+$, $d_v=p_V^- -p_V^+$,
\begin{equation}
 \begin{split}
 F_\theta&=R_Sd_r\cdot e_\theta
             +d_v\cdot(v_re_\theta-v_te_r),\\
 F_{v_r}&=d_v\cdot e_r,\qquad F_{v_t}=d_v\cdot e_\theta,\\
 F_\tau&=H^+-H^-+d_r\cdot V_p+d_v\cdot A_p.
 \end{split}
 \label{eq:soichart}
\end{equation}
At a regular optimized crossing these derivatives vanish.
The position jump in \eqref{eq:soijump} is normal to the surface, as
required by the angle variation.

Translation changes the Hamiltonian to
\begin{equation}
 H_{\rm rel}=H_{\rm global}-p_X\cdot V_p-p_V\cdot A_p,
\end{equation}
which is continuous at the interface. For a given continuous control
the state variation also has saltation matrix
\begin{equation}
 \begin{gathered}
 S=I+\frac{(f^+-f^-)n^T}{n^Tf^-+c_{S,t}},\\
 n=(x,0,0),\qquad f^+-f^-=(0,\Delta G,0).
 \end{gathered}
\end{equation}
It reproduces the covector jump through $p^-=S^Tp^+$.
Full optimal-control shooting derivatives additionally need the
selector and adjoint variations.

\subsection{Units and mass across a route}
Let $\lambda=L/R_0$ and $\gamma_\mu=\mu_p/\mu_\odot$.
The local-to-Sun ratios are
\begin{equation}
 c_t=\frac{\lambda^{3/2}}{\sqrt{\gamma_\mu}},\quad
 c_v=\sqrt{\gamma_\mu/\lambda},\quad
 c_J=\frac{\gamma_\mu^{3/2}}{\lambda^{5/2}}.
\end{equation}
For actions normalized by their respective effort scales,
\begin{equation}
 p_X^{(0)}=\frac{c_J}{\lambda}p_x^{(p)},\quad
 p_V^{(0)}=\frac{c_J}{c_v}p_v^{(p)},\quad
 H^{(0)}=\frac{c_J}{c_t}H^{(p)}.
 \label{eq:covscale}
\end{equation}
The resource covector is unchanged when resource and action use the
same effort scale. Apply any frame rotation as well.
This follows by preserving the physical endpoint differential, not
by matching raw normalized covectors.

For equal power on every leg, total physical effort and time are sums
$\sum_i J_{0,i}J_i$ and $\sum_i T_{0,i}h_i$.
For unequal constant powers without mass resets,
\begin{equation}
 \frac1{m_f}-\frac1{m_0}
 =\sum_i\frac{J_{0,i}J_i}{2P_i}.
 \label{eq:route_mass}
\end{equation}
Each leg starts at the preceding leg's actual final mass.
Resource states cannot be reset to their original launch values.
A moon patch uses its prescribed origin acceleration and the actual
parent field outside; do not add the origin acceleration twice.

A prescribed stop, a point flyby and an SOI crossing have different
interface conditions. Once all interior variables are stationary,
the envelope theorem leaves only the outer endpoint terms in a route
timing derivative. It does not make each embedded phase satisfy $H=0$.
